\section{Ran(\(E\)) obstruction and compact/hybrid repair}
\label{sec:ranE-obstruction-attack-repair}

\subsection{Claim attacked}

The attacked claim is the naive implication
\[
\mathcal F_X\in\mathrm{Fact}(\operatorname{Ran}(E))
\quad\Longrightarrow\quad
\hbox{all Igusa charge directions are support-local on }E.
\]
This implication is false.  The projection-finite sector is local on
\(\operatorname{Ran}(E)\).  Positive elliptic degree is global over
\(E\), and ordinary support-local factorization does not see it.

\subsection{Exact obstruction}

\begin{theorem}[Elliptic-degree locality obstruction]
Let \(p_E:K3\times E\to E\) be projection.  Let
\(C\subset K3\times E\) be a one-dimensional subscheme whose curve class
has positive elliptic degree:
\[
(p_E)_*[C]=b[E],\qquad b>0.
\]
Then \(p_E(C)=E\).  Hence \(C\) is not supported over a finite subset of
\(E\), and it is not contained in \(p_E^{-1}(U)\) for any proper open
\(U\subsetneq E\).
\end{theorem}

\begin{proof}
If \(p_E(C)\) were zero-dimensional, the pushforward of the one-cycle
\([C]\) to the curve \(E\) would be zero.  This contradicts
\((p_E)_*[C]=b[E]\) with \(b>0\).  Thus \(p_E(C)\) is a one-dimensional
closed subset of the irreducible curve \(E\).  Therefore \(p_E(C)=E\).
\end{proof}

\subsection{Sector distinction}

A sector is projection-finite over \(E\) when every object \(A\) in it
satisfies
\[
p_E(\operatorname{Supp}A)\subset F
\]
for some finite \(F\subset E\).  This is the sector where disjoint-open
factorization on \(\operatorname{Ran}(E)\) has its usual support
meaning.

A sector is wrapped or global over \(E\) when the one-dimensional part of
the charge has positive elliptic degree.  These objects dominate \(E\).
They cannot be separated by disjoint opens of \(E\), so the usual Ran
diagonal calculus does not by itself produce their collision theory.

The \(s\)-variable of the Igusa product is the OP elliptic-degree
direction after the variable conversion in the manuscript.  Therefore a
naive projection-local Hall/chiral object misses the positive
\(s\)-direction unless an additional global mechanism is supplied.

\subsection{Repair alternatives}

There are two honest repair routes.

\begin{enumerate}
\item \textbf{Projection-finite local factorization plus global
Fock/Hecke modes.}  Construct the local Hall/chiral factorization object
only for the projection-finite sector.  Then construct a separate global
Fock/Hecke mechanism carrying positive elliptic degree, and prove its
compatibility with protected integration, the \(s\)-variable, and the
primitive denominator comparison.

\item \textbf{Hybrid local/global wrapped factorization.}  Enlarge the
factorization object.  Local vertical insertions live on
\(\operatorname{Ran}(E)\); positive elliptic-degree sectors are global
wrapped legs.  The missing theorem is the extension-stack construction
with local and global legs, together with Thom--Sebastiani orientation
transport, associativity over two-step flags, and compatibility with
protected integration.
\end{enumerate}

Both routes preserve the Borcherds determinant.  The obstruction is not a
defect in
\[
\Delta_5=\operatorname{Borch}_{\theta}(\phi_{0,1})
\]
or in the Gritsenko--Nikulin denominator identity.  It is a realization
obstruction for a compact \(E_3\)/Hall/chiral source.

\subsection{Main-text integration}

The manuscript now names the obstruction in Section
\ref{sec:physical-question}.  The same obstruction is stated formally in
the holomorphic \(E_3\)-to-chiral discussion before the compact Igusa
datum:
\[
\text{projection-finite Ran locality}
\quad\hbox{versus}\quad
\text{wrapped/global elliptic-degree sectors}.
\]

\subsection{Files changed}

\begin{itemize}
\item \verb|main.tex|: Section \ref{sec:physical-question}; the
holomorphic \(E_3\)/factorization discussion before
Definition~\ref{def:compact-igusa-realization-datum}.
\item \verb|agent_material/13_ranE_obstruction_attack_repair.tex|.
\item \verb|out/main.pdf|: rebuilt by \verb|make|.
\end{itemize}

\subsection{Checks run}

\begin{itemize}
\item Read \verb|CLAUDE.md|, \verb|AGENTS.md|,
\verb|/Users/raeez/ecosystem/INVARIANTS.md|, \verb|main.tex|,
\verb|proj.bib|, \verb|agent_material/12_compact_hall_ranE_attack_heal.tex|,
and \verb|agent_material/12_synthesis_architecture_attack_heal.tex|.
\item Read the harness self-reflection section in
\verb|/Users/raeez/ecosystem/AGENTS-HARNESS.md|.
\item Grepped \verb|main.tex| and \verb|agent_material| for
\(\operatorname{Ran}(E)\), elliptic degree, support locality, Fock,
Hecke, and compact realization anchors.
\item Ran \verb|git diff --check|; no whitespace errors were reported.
\item Ran \verb|make|; it completed successfully and wrote
\verb|out/main.pdf|.  The TeX log still contains layout warnings
\(overfull/underfull boxes and a hyperref PDF-string warning\), but no
fatal error.
\end{itemize}

\subsection{Remaining mathematical obligations}

\begin{enumerate}
\item Prove a projection-finite Hall/chiral locality theorem with the
chosen compact heart, stability condition, HN completion, critical atlas,
and orientation.
\item If the first repair route is chosen, construct the global
Fock/Hecke elliptic-degree modes and prove compatibility with protected
integration and the Borcherds product.
\item If the hybrid route is chosen, construct extension stacks with
local Ran legs and wrapped \(E\)-global legs, then prove associativity,
orientation compatibility, and continuity in the sectorial completion.
\item Prove that either repaired source has protected primitive
superdimensions \(f(nm,l)\) and primitive bracket
\(\mathfrak g_{\Delta_5}\), rather than importing these from the
denominator identity.
\end{enumerate}
